\clearpage
\fancyhead[L]{\small BEOBACHTUNGEN.}
\SourceLeaf{471}
\begin{center}
\rule{0.55\linewidth}{0.4pt}\\[2pt]
Astronomische Nachrichten. Band XXII. Nr. 506. Seite 31. 1844 Juni 20.\\[2pt]
\rule{0.55\linewidth}{0.4pt}
\end{center}
\begin{center}
Beobachtung der Mondfinsterniss vom 31. Mai 1844 zu Göttingen.\\
Die Zeiten sind in Sternzeit angesetzt.
\end{center}

\begin{center}\begingroup\fontsize{10}{12}\selectfont\setlength{\tabcolsep}{5pt}
\begin{tabular}{@{}p{0.47\linewidth}|r|r@{}}
\emph{Eintritt} & \multicolumn{1}{c|}{\textsc{Gauss}} & \multicolumn{1}{c}{\textsc{Goldschmidt}}\\\hline
Anfang\ \mbox{.\,.\,.\,.} & $14^{\rm u}\ 27^{\rm m}$ & \\
Riccioli\ \mbox{.\,.\,.\,.} & $32\ 24^{\rm s}$ & \\
Aristarch\ \mbox{.\,.\,.\,.} & $32\ 48$ & $14^{\rm u}\ 32^{\rm m}\ 36^{\rm s}$\\
Grimaldi\ \mbox{.\,.\,.\,.} & $34\ 4$ & \\
Kepler\ \mbox{.\,.\,.\,.} & $37\ 30$ & $37\ 15$\\
Mayer\ \mbox{.\,.\,.\,.} &  & $38\ 57$\\
Crüger\ \mbox{.\,.\,.\,.} & $39\ 2$ & \\
Pytheas\ \mbox{.\,.\,.\,.} & $42\ 1$ & $41\ 46$\\
Copernicus I.\ \mbox{.\,.} & $42\ 57$ & $42\ 50$\\
Plato I.\ \mbox{.\,.\,.\,.} & $44\ 3$ & $43\ 51$\\
Plato II.\ \mbox{.\,.\,.\,.} & $44\ 58$ & $45\ 18$\\
Copernicus II.\ \mbox{.\,.} &  & $45\ 5$\\
Gassendi I.\ \mbox{.\,.\,.} & $45\ 16$ & \\
Gassendi II.\ \mbox{.\,.\,.} & $46\ 6$ & \\
Calippus\ \mbox{.\,.\,.\,.} & $51\ 48$ & \\
Aristoteles\ \mbox{.\,.\,.} & $52\ 35$ & \\
Eudoxus\ \mbox{.\,.\,.\,.} & $53\ 12$ & \\
Sulpicius Gallus\ \mbox{.} & $54\ 38$ & \\
Manilius I.\ \mbox{.\,.\,.} & $55\ 36$ & $55\ 25$\\
Manilius II.\ \mbox{.\,.\,.} & $56\ 23$ & $56\ 49$\\
Menelaus\ \mbox{.\,.\,.} & $59\ 5$ & $58\ 39$\\
Posidonius I.\ \mbox{.\,.} & $59\ 54$ & \\
Posidonius II.\ \mbox{.\,.} & $15\ 1\ 7$ & \\
Plinius I.\ \mbox{.\,.\,.} &  & $15\ 1\ 21$\\
Plinius II.\ \mbox{.\,.\,.} &  & $2\ 24$\\
Dionysius\ \mbox{.\,.\,.} & $2\ 49$ & $-2\ 58$\\
Tycho I.\ \mbox{.\,.\,.\,.} & $6\ 2$ & $5\ 49$\\
Tycho II.\ \mbox{.\,.\,.} & $7\ 38$ & $7\ 50$\\
Censorinus\ \mbox{.\,.\,.} & $10\ 52$ & \\
Proclus\ \mbox{.\,.\,.\,.} & $11\ 23$ & $11\ 4$\\
Taruntius\ \mbox{.\,.\,.} & $14\ 29$ & \\
Goclenius I.\ \mbox{.\,.\,.} & $17\ 0$ & $16\ 56$\\
Goclenius II.\ \mbox{.\,.} & $17\ 36$ & $17\ 50$\\
Anfang der totalen Verfinsterung\ \mbox{.} & $29\ 28$ & $30\ 4$\\
\end{tabular}\endgroup\end{center}

\clearpage
\begin{center}\begingroup\fontsize{10}{12}\selectfont\setlength{\tabcolsep}{5pt}
\begin{tabular}{@{}p{0.47\linewidth}|r|r@{}}
\emph{Austritte.} & \multicolumn{1}{c|}{\textsc{Gauss}} & \multicolumn{1}{c}{\textsc{Goldschmidt}}\\\hline
Ende der totalen Verfinsterung\ \mbox{.} & $16^{\rm u}.46^{\rm m}\ 54^{\rm s}$ & \\
Crüger\ \mbox{.\,.\,.\,.} & $49\ 34$ & \\
Riccioli\ \mbox{.\,.\,.\,.} & $51\ 1$ & \\
Grimaldi I.\ \mbox{.\,.\,.} &  & $16^{\rm u}\ 50^{\rm m}\ 29$\\
Grimaldi II.\ \mbox{.\,,\,.} &  & $16\ 51\ 17$\\
Marius\ \mbox{.\,.\,.\,.} &  & $55\ 44$\\
Gassendi I.\ \mbox{.\,.\,.} & $57\ 14$ & \\
Gassendi II.\ \mbox{.\,.\,.} & $58\ 36$ & \\
Kepler I.\ \mbox{.\,.\,.\,.} &  & $17\ 0\ 22$\\
Kepler II.\ \mbox{.\,.\,.} &  & $1\ 58$\\
Aristarch\ \mbox{.\,.\,.} & $17\ 2\ 41$ & $2\ 22$\\
Tycho I.\ \mbox{.\,.\,.\,.} & $4\ 51$ & $3\ 59$\\
Tycho II.\ \mbox{.\,.\,.\,.} & $6\ 31$ & $6\ 12$\\
Mayer\ \mbox{.\,.\,.\,.\,.} &  & $8\ 27$\\
Copernicus I.\ \mbox{.\,.} & $10\ 20$ & $10\ 8$\\
Copernicus II.\ \mbox{.\,.} & $11\ 56$ & $11\ 55$\\
Pytheas\ \mbox{.\,.\,.\,.} & $12\ 33$ & $12\ 34$\\
Laplace\ \mbox{.\,.\,.\,.} & $14\ 56$ & \\
Plato I.\ \mbox{.\,.\,.\,.} & $19\ 53$ & $19\ 36$\\
Plato II.\ \mbox{.\,.\,.\,.} & $20\ 58$ & $20\ 41$\\
Manilius I.\ \mbox{.\,.\,.} &  & $25\ 38$\\
Manilius Mitte & $25\ 58$ & \\
Manilius II.\ \mbox{.\,.\,.} &  & $26\ 27$\\
Sulpicius Gallus\ \mbox{.} & $26\ 54$ & \\
Dionysius\ \mbox{.\,.\,.} & $28\ 49$ & $28\ 13$\\
Menelaus I.\ \mbox{.\,.\,.} &  & $28\ 46$\\
Menelaus Mitte\ \mbox{.} & $29\ 25$ & \\
Menelaus II.\ \mbox{.\,.} &  & $29\ 31$\\
Fracastos\ \mbox{.\,.\,.\,.} & $31\ 44$ & \\
Plinius\ \mbox{.\,.\,.\,.} & $32\ 59$ & $32\ 46$\\
Posidonius I.\ \mbox{.\,.} & $34\ 49$ & \\
Posidonius II.\ \mbox{.\,.} & $35\ 57$ & $35\ 58$\\
Censorinus\ \mbox{.\,.\,.} & $36\ 4$ & \\
Goclenius I.\ \mbox{.\,.} & $38\ 44$ & \\
Goclenius II.\ \mbox{.\,.} & $39\ 14$ & $40\ 6$\\
Jansen A (nach Mädler) & $39\ 53$ & \\
Proclus\ \mbox{.\,.\,.\,.} & $42\ 16$ & $41\ 59$\\
Picard\ \mbox{.\,.\,.\,.} & $44\ 53$ & \\
Ende\ \mbox{.\,.\,.\,.\,.} & $50$ & $49\ 24$\\
\end{tabular}\endgroup\end{center}

Während der totalen Verfinsterung erschien der vom Mittelpunkt des Schattens entfernte Mondrand in schönem lichten Blau, das übrige in dunklem Kupferroth.
\begin{flushright}\textsc{Gauss.}\end{flushright}
